\begin{ex}
Trong phép chia $12 : 2 = 6$, số 6 được gọi là gì? Tính thương của 12 và 2.
\shortans{6}
\loigiai{
Thương là 6.
}
\end{ex}

\begin{ex}
Trong phép chia $20 : 5 = 4$, số bị chia bằng bao nhiêu?
\shortans{20}
\loigiai{
Số bị chia là 20.
}
\end{ex}

\begin{ex}
Trong phép chia $35 : 5 = 7$, số chia bằng bao nhiêu?
\shortans{5}
\loigiai{
Số chia là 5.
}
\end{ex}

\begin{ex}
Biết số bị chia là 16, số chia là 2. Tính thương.
\shortans{8}
\loigiai{
Thương: $16 : 2 = 8$.
}
\end{ex}

\begin{ex}
Biết số bị chia là 40, số chia là 5. Tính thương.
\shortans{8}
\loigiai{
Thương: $40 : 5 = 8$.
}
\end{ex}

\begin{ex}
Biết số bị chia là 18, số chia là 2. Tính thương.
\shortans{9}
\loigiai{
Thương: $18 : 2 = 9$.
}
\end{ex}

\begin{ex}
Biết số bị chia là 45, số chia là 5. Tính thương.
\shortans{9}
\loigiai{
Thương: $45 : 5 = 9$.
}
\end{ex}

\begin{ex}
Biết thương là 7, số chia là 2. Tìm số bị chia.
\shortans{14}
\loigiai{
Số bị chia: $7 \times 2 = 14$.
}
\end{ex}

\begin{ex}
Biết thương là 6, số chia là 5. Tìm số bị chia.
\shortans{30}
\loigiai{
Số bị chia: $6 \times 5 = 30$.
}
\end{ex}

\begin{ex}
Biết thương là 9, số chia là 2. Tìm số bị chia.
\shortans{18}
\loigiai{
Số bị chia: $9 \times 2 = 18$.
}
\end{ex}

\begin{ex}
Biết thương là 8, số chia là 5. Tìm số bị chia.
\shortans{40}
\loigiai{
Số bị chia: $8 \times 5 = 40$.
}
\end{ex}

\begin{ex}
Biết số bị chia là 10, thương là 5. Tìm số chia.
\shortans{2}
\loigiai{
Số chia: $10 : 5 = 2$.
}
\end{ex}

\begin{ex}
Biết số bị chia là 25, thương là 5. Tìm số chia.
\shortans{5}
\loigiai{
Số chia: $25 : 5 = 5$.
}
\end{ex}

\begin{ex}
Tính giá trị biểu thức: $16 : 2 + 15$.
\shortans{23}
\loigiai{
Ta có $16 : 2 = 8$, rồi $8 + 15 = 23$.
}
\end{ex}

\begin{ex}
Tính giá trị biểu thức: $30 : 5 + 24$.
\shortans{30}
\loigiai{
Ta có $30 : 5 = 6$, rồi $6 + 24 = 30$.
}
\end{ex}

\begin{ex}
Tính giá trị biểu thức: $14 : 2 \times 3$.
\shortans{21}
\loigiai{
Ta có $14 : 2 = 7$, rồi $7 \times 3 = 21$.
}
\end{ex}

\begin{ex}
Tính giá trị biểu thức: $40 : 5 \times 2$.
\shortans{16}
\loigiai{
Ta có $40 : 5 = 8$, rồi $8 \times 2 = 16$.
}
\end{ex}

\begin{ex}
Điền số thích hợp: $18 : \dots = 9$.
\shortans{2}
\loigiai{
Số cần điền là: $18 : 9 = 2$.
}
\end{ex}

\begin{ex}
Điền số thích hợp: $35 : \dots = 7$.
\shortans{5}
\loigiai{
Số cần điền là: $35 : 7 = 5$.
}
\end{ex}

\begin{ex}
Điền số thích hợp: $\dots : 2 = 8$.
\shortans{16}
\loigiai{
Số cần điền là: $8 \times 2 = 16$.
}
\end{ex}

\begin{ex}
Điền số thích hợp: $\dots : 5 = 9$.
\shortans{45}
\loigiai{
Số cần điền là: $9 \times 5 = 45$.
}
\end{ex}

\begin{ex}
Trong một phép chia có số chia là 2, thương là số lớn nhất có một chữ số. Số bị chia là bao nhiêu?
\shortans{18}
\loigiai{
Số lớn nhất có 1 chữ số là 9. Số bị chia: $9 \times 2 = 18$.
}
\end{ex}

\begin{ex}
Trong một phép chia có số chia là 5, thương là số chẵn lớn nhất có một chữ số. Số bị chia là bao nhiêu?
\shortans{40}
\loigiai{
Số chẵn lớn nhất có 1 chữ số là 8. Số bị chia: $8 \times 5 = 40$.
}
\end{ex}

\begin{ex}
Một số chia cho 2 thì được 10. Lấy số đó trừ đi 5 thì được bao nhiêu?
\shortans{15}
\loigiai{
Số đó là: $10 \times 2 = 20$. Hiệu: $20 - 5 = 15$.
}
\end{ex}

\begin{ex}
Tính: $50 : 5 - 2$.
\shortans{8}
\loigiai{
Ta có $50 : 5 = 10$, rồi $10 - 2 = 8$.
}
\end{ex}
